\chapter{Топологія \texorpdfstring{$\psi$}{ψ}: O- і X-точки}\label{ch:g05}

GS-нев'язка з розд.~\ref{ch:g04} спрацьовує вже на 1\,\% шуму, але вона кількісна. Цей розділ про якісний, бінарний критерій: чи має подане поле взагалі \emph{топологію} рівноваги токамака? Тут розглянуто пункти E9--E13 (з уточненням N4, внесеним у реєстр 2026-09-29), спростування R5 і перевірку V-D4. Геометрію O- і X-точок див. у~\cite{Manual}, розд.~4.

\section{Фізика: критичні точки, індекс, Пуанкаре--Хопф, Морс}\label{sec:g05-physics}

Критична точка $\psi(R,Z)$~--- це точка, де $\grad\psi=0$, тобто полоїдальне поле $\Bpol=|\grad\psi|/R$ дорівнює нулю. Тип точки визначає гесіан $\mathsf H_\psi$. Якщо $\det\mathsf H_\psi>0$, маємо екстремум: це \emph{O-точка} (еліптична). Навколо неї лежать вкладені магнітні поверхні. Магнітна вісь~--- саме така точка, і зайва O-точка означає магнітний острів. Якщо $\det\mathsf H_\psi<0$, маємо сідло: це \emph{X-точка} (гіперболічна), через яку проходить сепаратриса.

\paragraph{Індекс (winding number).} Обійдемо малий контур навколо ізольованої критичної точки і порахуємо, скільки обертів робить вектор $\grad\psi$. Для екстремуму індекс дорівнює $+1$, причому \emph{однаково} для максимуму й мінімуму (при $\psi\to-\psi$ вектор просто повертається на $\pi$). Для сідла індекс $-1$. Саме тому критерій не залежить від знакової конвенції машини, що знадобиться в розд.~\ref{ch:g06}.

\begin{outofcorpus}[теорема Пуанкаре--Хопфа і теорія Морса]
\emph{Пуанкаре--Хопф.} Нехай поле $\grad\psi$ на межі області $D$ всюди напрямлене назовні (або всюди всередину) і має скінченну кількість нулів. Тоді сума індексів нулів у $D$ дорівнює ейлеровій характеристиці $\chi_{\mathrm E}(D)$. Для диска $\chi_{\mathrm E}=1$, тож усередині замкненої магнітної поверхні $N_O-N_X=1$.

\emph{Морс.} Для функції Морса ейлерова характеристика підрівневої множини $\{\psi\le t\}$ змінюється лише тоді, коли $t$ проходить критичне значення: на $+1$ при мінімумі чи максимумі (індекс Морса 0 чи 2) і на $-1$ при сідлі (індекс Морса 1). Отже,
\begin{equation}\label{eq:g05-morse}
  \chi_{\mathrm E}\big(\{\psi\le t\}\cap D\big)=\sum_{c:\ \psi(c)\le t}(-1)^{\lambda_c}.
\end{equation}
Тобто $\chi_{\mathrm E}$-крива~--- це \emph{інтеграл} (кумулятивна сума) лічильника критичних точок зі знаком.
\end{outofcorpus}

Звідси три наслідки, які нижче перевірено числами.
\begin{enumerate}
\item У справжній рівновазі всередині LCFS маємо рівно одну O-точку і жодного сідла, сума індексів $+1$ (E9).
\item Шум народжує критичні точки \emph{парами}: у складці (saddle-node) з'являються екстремум і сідло разом. Кількість зростає, а сума індексів $+1-1=0$ не змінюється. Отже, сума як метрика марна, інформацію несе лише кількість (E11).
\item Пара O--X, що народилася поруч, дає в \eqref{eq:g05-morse} внески $+1$ і $-1$, які майже одразу скорочуються. Тому $\chi_{\mathrm E}$-крива менш чутлива за лічильник (E13, R5).
\end{enumerate}

\section{Звідки ідея}\label{sec:g05-origin}

Ідея прийшла з аналізу перерізу Пуанкаре (розд.~\ref{ch:g07}). В осесиметрії система силових ліній інтегровна, і переріз дає рівно контури $\psi$ (E18). Тому динамічний зміст, що лишається у 2D,~--- це не сам переріз, а \emph{скелет критичних точок} $\psi$. Так мотивовано docstring \texttt{topology\_probe.py}.

\emph{Увага:} до 2026-09-29 у рядках 6--7 цього docstring стояло «psi as canonical momentum»~--- та сама помилка, що R4 (правильно: імпульс~--- $\psit$, гамільтоніан~--- $\psi$, розд.~\ref{ch:g07}). На код вона не впливала; docstring виправлено 2026-09-29 без зміни кількості рядків~\cite{RegLog}.

Новизну перевірено в журналі (V-D2~\cite{RegVerif}): метод опубліковано в сегментації (Clough та ін., топологічна втрата без міток, лише з апріорної топології) і в стисненні даних (cpSZ). Нове лише \emph{застосування}: ніхто не робить кількість зайвих критичних точок метрикою фузійних сурогатів, і жоден фузійний бенчмарк не має топологічного члена.

\section{Де це в коді}\label{sec:g05-code}

\code{ourtry/04-novelty/topology_probe.py}{39}{48}{(індекс як число обертів $\grad\psi$)}

Рядки 41--44 обходять периметр прямокутника $[i_{z0},i_{z1}]\times[i_{r0},i_{r1}]$ проти годинникової стрілки. Рядок 45 бере кут $\arg\grad\psi$ у кожній точці периметра. Рядки 46--47 рахують прирости кута, зведені до $(-\pi,\pi]$, і це ключовий крок: без зведення стрибок $2\pi$ на розрізі $\arctan$ зіпсував би суму. Рядок 48 дає повне число обертів~--- індекс.

\code{ourtry/04-novelty/topology_probe.py}{51}{79}{(\texttt{critical\_points}: попіксельне детектування)}
\code{ourtry/04-novelty/topology_probe.py}{80}{90}{(продовження: зв'язні компоненти)}

\begin{itemize}
\item \textbf{Рядки 64--66}: $\grad\psi$ через \texttt{np.gradient} і кут \texttt{arctan2}.
\item \textbf{Рядки 69--78}: для кожного пікселя всередині області~--- число обертів по кільцю з восьми сусідів (\texttt{LOOP}, рядок 36). Ненульове значення означає, що в комірці сидить критична точка.
\item \textbf{Рядки 80--89}~--- \emph{виправлення V-D4}. Перша версія детектора рахувала попіксельні детекції як окремі точки і на істині давала \textbf{4 точки замість 1}. Діагностика показала кластер $2\times2$ навколо \emph{тієї самої} магнітної осі: вісь ($R=1{,}678$, $Z=+0{,}030$) лежить між пікселями, тож кільця чотирьох сусідніх пікселів усі її охоплюють~\cite{RegVerif}. Виправлення: детекції об'єднуються у зв'язні компоненти (\texttt{scipy.ndimage.label} зі зв'язністю 8), і для кожної компоненти рахується число обертів по периметру її обмежувального прямокутника, розширеного на одну комірку. Це математично коректний індекс усього, що лежить усередині.
\end{itemize}

\code{ourtry/04-novelty/topology_probe.py}{93}{102}{(область аналізу)}

Область~--- внутрішність LCFS, видобутої процедурою скорера, з ерозією на дві комірки (рядок 102). Ерозія тут не косметика. Вона виключає X-точку, що лежить \emph{на} сепаратрисі, і відсуває межу області всередину, де $\grad\psi\neq0$ і (для гладкого $\psi$) трансверсальний до межі~--- саме цього вимагає Пуанкаре--Хопф. Строго для піксельної області й зашумленого поля це не доведено, лише правдоподібно (коментар путівника). Журнал V-D4 окремо наголошує: сума $+1$ виходить \emph{тому}, що X-точки виключено. На прямокутнику, що містить $k$ X-точок, сума дорівнює $1-k$, і в конфігурації double-null критерій без явно зафіксованої області ламається.

\section{Що вийшло: E9--E11}\label{sec:g05-results}

\begin{established}{E9, E10, E11}
\textbf{E9.} Скелет істини точно канонічний: рівно 1 еліптична точка, 0 сідел, сума індексів $+1$.
\textbf{E10.} Скелет стійкий до 3\,\% шуму і руйнується на 10\,\%: при $R^2_\psi=0{,}993$ маємо 20,5 зайвих нерухомих точок.
\textbf{E11.} Сума індексів зберігається навіть у хаосі. Пуанкаре--Хопф змушує точки народжуватися парами O--X, тож як метрика сума марна, працює лише кількість~\cite{RegEst}.
\end{established}

Повторний запуск 2026-09-29 (\texttt{topology\_probe.py}, 8 кадрів розряду 203702, зерна $1000+k$) відтворив E9--E11. До рівня шуму 3\,\% включно медіани дорівнюють 1~O, 0~X, сума $1{,}0$ і 0 зайвих. На 10\,\% ($R^2_\psi=0{,}99278$) маємо 10,5~O, 11,0~X, медіану суми 1,5 і 20,5 зайвих. Рис.~\ref{fig:g05-cp} показує один кадр.

\begin{figure}[htbp]\centering
\includegraphics[width=0.8\textwidth]{g05_critical_points.jpg}
\caption[Критичні точки psi: істина і 10 \% шуму]{Критичні точки $\psi$ усередині LCFS (синій контур~--- область після ерозії) для кадру 75 розряду DIII-D 203702. Ліворуч: EFIT, рівно одна O-точка (магнітна вісь). Праворуч: +10\,\% білого шуму, 8~O ($\bullet$) і 7~X ($\times$), сума індексів $+1$. Сірим показано контури істинного $\psi$. Власний розрахунок функціями \texttt{inside\_lcfs}, \texttt{critical\_points} з \texttt{topology\_probe.py}; рисунок~--- \texttt{guide/figures/g05\_fig.py}. Дані: Fusion Equilibrium Challenge (Sophelio / General Atomics), CC BY 4.0.}\label{fig:g05-cp}
\end{figure}

\paragraph{Уточнення N4 (внесено в реєстр 2026-09-29).} Поріг 3\,\% \emph{пограничний}. \texttt{topology\_probe.py} на 3\,\% дає медіану 1~O/0~X, тобто лічильник мовчить. Натомість \texttt{topology\_ec\_curve.py} на іншій реалізації шуму (зерна $7000+k$, 6 кадрів) дає медіану лічильника 1,5, тобто лічильник спрацьовує. Обидва скрипти перезапущено, числа відтворились. Тому «стійкий до 3\,\%» слід читати як «на межі стійкості». Для E11 медіана суми при 10\,\% дорівнює 1,5, тобто «близька до $+1$», а не рівна $+1$. Висновок «сума марна, працює кількість» від цього не змінюється~\cite{RegEst}. Одне спостереження з нашого прогону: на 10\,\% другий скрипт дає медіану 17,0 критичних точок, тобто 16 зайвих, проти 20,5 у першого. Масштаб той самий, але точне число залежить від реалізації шуму й набору кадрів.

\section{Ейлерова характеристика: еталон без налаштування (E12) і спростування R5}\label{sec:g05-ec}

\code{ourtry/04-novelty/topology_ec_curve.py}{25}{32}{($\chi_{\mathrm E}$-крива підрівневих множин)}

Для кожного з 40 порогів $t$ маска $\{\psi\le t\}\cap D$ передається в \texttt{skimage.measure.euler\_number} зі зв'язністю 1. Складність $\order{N}$ на рівень. На DIII-D вісь є \emph{мінімумом} $\psi$ (\texttt{AXIS\_SIGN}$=-1$, розд.~\ref{ch:g06}), тож $\{\psi\le t\}$ поблизу осі~--- диск.

\code{ourtry/04-novelty/topology_ec_curve.py}{59}{85}{(порівняння на одній драбині шуму)}

Рядки 66--69: $\chi_{\mathrm E}$-криву збуреного поля рахують на \emph{еталонній} масці й еталонних порогах, щоб метрику не збивало зміщення передбаченої LCFS. Рядки 70--74 обчислюють середнє і максимальне $|\Delta\chi_{\mathrm E}|$ та кількість критичних точок тим самим детектором. Рядки 77--84 формулюють вердикт: «EC fires», якщо середнє $|\Delta\chi_{\mathrm E}|>1$; «count fires, EC quiet», якщо критичних точок більше однієї.

\begin{established}{E12, E13}
\textbf{E12.} На істині $\chi_{\mathrm E}(\psi\le t)\equiv1$ на всіх рівнях: кожна підрівнева множина є диском. Це еталон без налаштування, для якого не треба знати кількість O-точок і не потрібен поріг персистентності.
\textbf{E13.} $\chi_{\mathrm E}$-крива менш чутлива за лічильник: вона мовчить на 3\,\%, де лічильник спрацьовує. За Морсом вона є його інтегралом~\cite{RegEst}.
\end{established}

Прогін 2026-09-29: на 3\,\% ($R^2_\psi=0{,}99956$) середнє $|\Delta\chi_{\mathrm E}|$ дорівнює 0,138 при медіані лічильника 1,5 (вердикт «count fires, EC quiet»). На 10\,\% маємо 4,075, тобто спрацьовують обидва. На істині $\chi_{\mathrm E}=1$ у всіх надрукованих рівнях.

\begin{refuted}{R5}
«$\chi_{\mathrm E}$-крива краща за лічильник критичних точок». Це була наша власна рекомендація; журнал V-D4 фіксує, що вона хибна. Вимір показав протилежне: крива \emph{менш} чутлива, бо за \eqref{eq:g05-morse} вона є інтегральною версією лічильника, а не незалежним виміром~\cite{RegEst,RegVerif}. Вона не марна: $\chi_{\mathrm E}\equiv1$ зручно подавати як \emph{шлюз} (див. \texttt{metric-design.md}, розділ «Шлюз»). Діагностику ж, тобто скільки зайвих островів і де саме, дає лічильник.
\end{refuted}

Чутливість топологічного критерію нижча за GS-нев'язку: 3\,\% проти 1\,\% шуму. Звідси його роль у \texttt{metric-design.md}: не градуйований член, а \emph{критерій допустимості}, пізніший, але якісно інший сигнал. Він ловить острови, яких не існує.

\begin{practicum}[топологічний шлюз]
\textbf{Відтворити} (з \texttt{fusion equilibrium challenge/starter}, $\approx$40~с і $\approx$60~с): \texttt{.venv/bin/python "../../Our try/04-novelty/topology\_probe.py"} і \texttt{... topology\_ec\_curve.py}.

\textbf{Задача~1 (N4 кількісно).} Замініть одне зерно на кадр (рядок 126, \texttt{1000 + k}) циклом по 20 зернах і оцініть \emph{імовірність} того, що лічильник спрацює на 1, 2, 3 і 5\,\% шуму. Поріг стане кривою ймовірності замість «так/ні».

\textbf{Задача~2 (персистентність).} \texttt{metric-design.md} вимагає персистентного порогу замість склеювання компонент, інакше детектор галюцинує точки в SOL. Відкидайте пари O--X, для яких $|\psi_O-\psi_X|$ менший за поріг $\tau$. За якого $\tau$ істина лишається $1/0$, а 10\,\% шуму ще ловиться?

\textbf{Задача~3 (double-null).} Візьміть прямокутник, що містить X-точку, без ерозії LCFS і переконайтеся, що сума індексів дорівнює $1-k$. Запропонуйте явне означення області для конфігурації з двома X-точками.
\end{practicum}
